\documentclass[10pt,a4paper]{amsart}
\usepackage{amsmath,amssymb,amsthm,graphicx}
\usepackage{verbatim}
\newenvironment{reference}{\comment}{\endcomment}
%\newenvironment{reference}{}{}
\newenvironment{slogan}{\comment}{\endcomment}
\newenvironment{history}{\comment}{\endcomment}

% For commutative diagrams we use Xy-pic
\usepackage[all]{xy}

% We use 2cell for 2-commutative diagrams.
\xyoption{2cell}
% Theorem environments.
%
\theoremstyle{plain}
\newtheorem{theorem}[subsection]{Theorem}
\newtheorem{proposition}[subsection]{Proposition}
\newtheorem{lemma}[subsection]{Lemma}

\theoremstyle{definition}
\newtheorem{definition}[subsection]{Definition}
\newtheorem{example}[subsection]{Example}
\newtheorem{exercise}[subsection]{Exercise}
\newtheorem{situation}[subsection]{Situation}

\theoremstyle{remark}
\newtheorem{remark}[subsection]{Remark}
\newtheorem{remarks}[subsection]{Remarks}

\numberwithin{equation}{subsection}

% Macros
%
\def\lim{\mathop{\mathrm{lim}}\nolimits}
\def\colim{\mathop{\mathrm{colim}}\nolimits}
\def\Spec{\mathop{\mathrm{Spec}}}
\def\Hom{\mathop{\mathrm{Hom}}\nolimits}
\def\Ext{\mathop{\mathrm{Ext}}\nolimits}
\def\SheafHom{\mathop{\mathcal{H}\!\mathit{om}}\nolimits}
\def\SheafExt{\mathop{\mathcal{E}\!\mathit{xt}}\nolimits}
\def\Sch{\mathit{Sch}}
\def\Mor{\mathop{\mathrm{Mor}}\nolimits}
\def\Ob{\mathop{\mathrm{Ob}}\nolimits}
\def\Sh{\mathop{\mathit{Sh}}\nolimits}
\def\NL{\mathop{N\!L}\nolimits}
\def\CH{\mathop{\mathrm{CH}}\nolimits}
\def\proetale{{pro\text{-}\acute{e}tale}}
\def\etale{{\acute{e}tale}}
\def\QCoh{\mathit{QCoh}}
\def\Ker{\mathop{\mathrm{Ker}}}
\def\Im{\mathop{\mathrm{Im}}}
\def\Coker{\mathop{\mathrm{Coker}}}
\def\Coim{\mathop{\mathrm{Coim}}}

% Boxtimes
%
\DeclareMathSymbol{\boxtimes}{\mathbin}{AMSa}{"02}

%
% Macros for moduli stacks/spaces
%
\def\QCohstack{\mathcal{QC}\!\mathit{oh}}
\def\Cohstack{\mathcal{C}\!\mathit{oh}}
\def\Spacesstack{\mathcal{S}\!\mathit{paces}}
\def\Quotfunctor{\mathrm{Quot}}
\def\Hilbfunctor{\mathrm{Hilb}}
\def\Curvesstack{\mathcal{C}\!\mathit{urves}}
\def\Polarizedstack{\mathcal{P}\!\mathit{olarized}}
\def\Complexesstack{\mathcal{C}\!\mathit{omplexes}}
% \Pic is the operator that assigns to X its picard group, usage \Pic(X)
% \Picardstack_{X/B} denotes the Picard stack of X over B
% \Picardfunctor_{X/B} denotes the Picard functor of X over B
\def\Pic{\mathop{\mathrm{Pic}}\nolimits}
\def\Picardstack{\mathcal{P}\!\mathit{ic}}
\def\Picardfunctor{\mathrm{Pic}}
\def\Deformationcategory{\mathcal{D}\!\mathit{ef}}
\usepackage{cleveref}
\crefname{originallemma}{original lemma}{original lemmas}
\Crefname{originallemma}{Original lemma}{Original lemmas}
\makeatletter
\newlabel{algebra-lemma-smooth-at-generic-point}{{07ND}{source}{}{sourcepointer-1}{}}
\newlabel{algebra-lemma-smooth-at-generic-point@cref}{{[originallemma][0][]07ND}{[1][0][]source}{}{}{}}
\makeatother
\crefname{originalreference}{original reference}{original references}
\Crefname{originalreference}{Original reference}{Original references}
\makeatletter
\newlabel{lemma-flat}{{00HD}{source}{}{sourcepointer-2}{}}
\newlabel{lemma-flat@cref}{{[originalreference][0][]00HD}{[1][0][]source}{}{}{}}
\makeatother
\crefname{originallemma}{original lemma}{original lemmas}
\Crefname{originallemma}{Original lemma}{Original lemmas}
\makeatletter
\newlabel{morphisms-lemma-flat-unramified-etale}{{02GV}{source}{}{sourcepointer-3}{}}
\newlabel{morphisms-lemma-flat-unramified-etale@cref}{{[originallemma][0][]02GV}{[1][0][]source}{}{}{}}
\makeatother
\crefname{originallemma}{original lemma}{original lemmas}
\Crefname{originallemma}{Original lemma}{Original lemmas}
\makeatletter
\newlabel{morphisms-lemma-smooth-omega-finite-locally-free}{{02G1}{source}{}{sourcepointer-4}{}}
\newlabel{morphisms-lemma-smooth-omega-finite-locally-free@cref}{{[originallemma][0][]02G1}{[1][0][]source}{}{}{}}
\makeatother
\crefname{originallemma}{original lemma}{original lemmas}
\Crefname{originallemma}{Original lemma}{Original lemmas}
\makeatletter
\newlabel{morphisms-lemma-triangle-differentials}{{01UX}{source}{}{sourcepointer-5}{}}
\newlabel{morphisms-lemma-triangle-differentials@cref}{{[originallemma][0][]01UX}{[1][0][]source}{}{}{}}
\makeatother
\crefname{originallemma}{original lemma}{original lemmas}
\Crefname{originallemma}{Original lemma}{Original lemmas}
\makeatletter
\newlabel{morphisms-lemma-unramified-omega-zero}{{02G5}{source}{}{sourcepointer-6}{}}
\newlabel{morphisms-lemma-unramified-omega-zero@cref}{{[originallemma][0][]02G5}{[1][0][]source}{}{}{}}
\makeatother
\title{Five equivalent conditions for a morphism of smooth curves to be separable, i.e. etale almost everywhere}
\author{The Stacks Project Authors}
\date{stacks-project:curves.tex:local-snapshot-20261012; GNU FDL 1.2 or later}
\begin{document}
\maketitle
\noindent\textit{Editorial note.} All delivered mathematical text is the author's own, copied from the pinned Stacks Project chapter file curves.tex; no statement or proof step was rewritten, reordered or removed.\par\medskip
\noindent\textit{Editorial note.} The counted proof is the author's complete argument: the differential triangle, the unramified and etale equivalences with automatic flatness, and the function-field criterion with the author's named citation and marked final detail are all retained verbatim.\par\medskip
\noindent\textit{Editorial note.} The label lemma-generically-etale names the counted statement, and the proof environment immediately following it in the source is delivered verbatim as its proof.\par\medskip
\noindent\textit{Editorial note.} The counted statement is carried by the author's own lemma environment in the chapter source.\par\medskip
\noindent\textit{Allowed standard prerequisites (editorial).}\begin{itemize}
\item Differentials on smooth curves. Morphisms, Lemma morphisms-lemma-smooth-omega-finite-locally-free and morphisms-lemma-triangle-differentials. Makes Omega\_\{X/k\} and Omega\_\{Y/k\} invertible and supplies the triangle f\textasciicircum{}*Omega\_\{Y/k\} -> Omega\_\{X/k\} -> Omega\_\{X/Y\} -> 0 relating conditions (1) and (2).
\item Unramified and etale loci. Morphisms, Lemmas morphisms-lemma-unramified-omega-zero and morphisms-lemma-flat-unramified-etale; Same chapter, Lemma lemma-flat. Equates the support condition with unramifiedness on an open and, flatness being automatic for curves, with etaleness on an open.
\item Function fields and generic smoothness. Algebra, Lemma algebra-lemma-smooth-at-generic-point. Identifies the etale-almost-everywhere condition with separability of the function field extension k(X)/k(Y).
\end{itemize}\medskip
\ifdefined\hypertarget\hypertarget{sourcepointer-1}{}\fi
\noindent\textit{Source reference.} Original lemma 07ND: sources/stacks-chow/tags.txt, source lines 9758--9758. Prerequisite pointer invoked at proof line(s) 2601; it names a standard fact of a sibling chapter that the counted argument uses, and every step of the counted proof built on it is retained verbatim in the delivered body.\par\medskip
\ifdefined\hypertarget\hypertarget{sourcepointer-2}{}\fi
\noindent\textit{Source reference.} Original reference 00HD: sources/stacks-curves/pinned-original-source.tex, source lines 303--303. Prerequisite pointer invoked at proof line(s) 2599; it names a standard fact of the same chapter that the counted argument uses, and every step of the counted proof built on it is retained verbatim in the delivered body.\par\medskip
\ifdefined\hypertarget\hypertarget{sourcepointer-3}{}\fi
\noindent\textit{Source reference.} Original lemma 02GV: sources/stacks-chow/tags.txt, source lines 3242--3242. Prerequisite pointer invoked at proof line(s) 2597; it names a standard fact of a sibling chapter that the counted argument uses, and every step of the counted proof built on it is retained verbatim in the delivered body.\par\medskip
\ifdefined\hypertarget\hypertarget{sourcepointer-4}{}\fi
\noindent\textit{Source reference.} Original lemma 02G1: sources/stacks-chow/tags.txt, source lines 3212--3212. Prerequisite pointer invoked at proof line(s) 2585; it names a standard fact of a sibling chapter that the counted argument uses, and every step of the counted proof built on it is retained verbatim in the delivered body.\par\medskip
\ifdefined\hypertarget\hypertarget{sourcepointer-5}{}\fi
\noindent\textit{Source reference.} Original lemma 01UX: sources/stacks-chow/tags.txt, source lines 2426--2426. Prerequisite pointer invoked at proof line(s) 2591; it names a standard fact of a sibling chapter that the counted argument uses, and every step of the counted proof built on it is retained verbatim in the delivered body.\par\medskip
\ifdefined\hypertarget\hypertarget{sourcepointer-6}{}\fi
\noindent\textit{Source reference.} Original lemma 02G5: sources/stacks-chow/tags.txt, source lines 3216--3216. Prerequisite pointer invoked at proof line(s) 2595; it names a standard fact of a sibling chapter that the counted argument uses, and every step of the counted proof built on it is retained verbatim in the delivered body.\par\medskip
\begin{lemma}
\label{lemma-generically-etale}
\begin{slogan}
A morphism of smooth curves is separable iff it is etale almost everywhere
\end{slogan}
Let $k$ be a field. Let $f : X \to Y$ be a morphism of smooth curves over $k$.
The following are equivalent
\begin{enumerate}
\item $\text{d}f : f^*\Omega_{Y/k} \to \Omega_{X/k}$ is nonzero,
\item $\Omega_{X/Y}$ is supported on a proper closed subset of $X$,
\item there exists a nonempty open $U \subset X$ such that
$f|_U : U \to Y$ is unramified,
\item there exists a nonempty open $U \subset X$ such that
$f|_U : U \to Y$ is \'etale,
\item the extension $k(X)/k(Y)$ of function fields is
finite separable.
\end{enumerate}
\end{lemma}

\begin{proof}
Since $X$ and $Y$ are smooth, the sheaves $\Omega_{X/k}$ and
$\Omega_{Y/k}$ are invertible modules, see
Morphisms, Lemma \ref{morphisms-lemma-smooth-omega-finite-locally-free}.
Using the exact sequence
$$
f^*\Omega_{Y/k} \longrightarrow \Omega_{X/k}
\longrightarrow \Omega_{X/Y} \longrightarrow 0
$$
of Morphisms, Lemma \ref{morphisms-lemma-triangle-differentials}
we see that (1) and (2) are equivalent and equivalent to the
condition that $f^*\Omega_{Y/k} \to \Omega_{X/k}$ is nonzero
in the generic point. The equivalence of (2) and (3) follows
from Morphisms, Lemma \ref{morphisms-lemma-unramified-omega-zero}.
The equivalence between (3) and (4) follows from
Morphisms, Lemma \ref{morphisms-lemma-flat-unramified-etale}
and the fact that flatness is automatic
(Lemma \ref{lemma-flat}).
To see the equivalence of (5) and (4)
use Algebra, Lemma \ref{algebra-lemma-smooth-at-generic-point}.
Some details omitted.
\end{proof}

\end{document}
